\documentclass{article}
\usepackage[T1]{fontenc}
\usepackage{lmodern}
\usepackage{graphicx}
\usepackage{amsmath,amssymb}
\usepackage[a4paper,margin=2cm]{geometry}
\usepackage[absolute,overlay]{textpos}
\setlength{\TPHorizModule}{1mm}
\setlength{\TPVertModule}{1mm}
\usepackage[indonesian]{babel}
\usepackage{hyperref}
\setlength{\emergencystretch}{2em}
% unit: Halaman 15

\begin{document}

% === Halaman 15 (python/native) ===

Pembentukan Kunci Putaran

• Kunci putaran ada sebanyak 2r + 2 buah yang masing-masing 

disimpan di dalam elemen-elemen larik yang dilabeli sebagai 

S[0], S[1], …, S[t – 1] dengan t = 2r + 2. 

• Setiap elemen larik panjangnya satu word (1 word = w bit)

15


Parameter 

Simbol 

Nilai yang dibolehkan 

Ukuran blok (dalam bit) 

w 

16, 32, 64 

Jumlah putaran 

r 

0, 1, …, 255 

Panjang kunci eksternal K 

(dalam byte, 1 byte = 8 bit) 

b 

0, 1, …, 255

\end{document}
